\documentclass[10pt,a4paper]{amsart}
\usepackage{amsmath,amssymb,amsthm,graphicx}
% Retained preamble for the Stacks Project chapter {CH} (shared project preamble).
% The chapter file pinned-original-source.tex only carries \input{{preamble}} on its first line,
% so no preamble text lives inside the chapter itself.  The Stacks Project uses one single shared
% preamble for every chapter: the snapshot stages it three times (sources/stacks-chow/preamble.tex,
% sources/stacks-etale-cohomology/preamble.tex and focused-seeds/AGP/sources/stacks-snapshot-20261011/
% preamble.tex) and those copies are identical in size.  The block below is copied mechanically,
% line for line, from sources/stacks-chow/preamble.tex lines 14-18 and 154-217: the Xy-pic setup and
% the native theorem declarations together with the mathematical macros.  Nothing was added, removed
% or redefined; the wrapper class and the amsmath/amssymb/amsthm/graphicx base packages are supplied
% by the assembler.
% For commutative diagrams we use Xy-pic
% Chapter: stacks-spaces-more-morphisms.tex -- More on Morphisms of Spaces.
\usepackage[all]{xy}

% We use 2cell for 2-commutative diagrams.
\xyoption{2cell}

% Theorem environments.
%
\theoremstyle{plain}
\newtheorem{theorem}[subsection]{Theorem}
\newtheorem{proposition}[subsection]{Proposition}
\newtheorem{lemma}[subsection]{Lemma}

\theoremstyle{definition}
\newtheorem{definition}[subsection]{Definition}
\newtheorem{example}[subsection]{Example}
\newtheorem{exercise}[subsection]{Exercise}
\newtheorem{situation}[subsection]{Situation}

\theoremstyle{remark}
\newtheorem{remark}[subsection]{Remark}
\newtheorem{remarks}[subsection]{Remarks}

\numberwithin{equation}{subsection}

% Macros
%
\def\lim{\mathop{\mathrm{lim}}\nolimits}
\def\colim{\mathop{\mathrm{colim}}\nolimits}
\def\Spec{\mathop{\mathrm{Spec}}}
\def\Hom{\mathop{\mathrm{Hom}}\nolimits}
\def\Ext{\mathop{\mathrm{Ext}}\nolimits}
\def\SheafHom{\mathop{\mathcal{H}\!\mathit{om}}\nolimits}
\def\SheafExt{\mathop{\mathcal{E}\!\mathit{xt}}\nolimits}
\def\Sch{\mathit{Sch}}
\def\Mor{\mathop{\mathrm{Mor}}\nolimits}
\def\Ob{\mathop{\mathrm{Ob}}\nolimits}
\def\Sh{\mathop{\mathit{Sh}}\nolimits}
\def\NL{\mathop{N\!L}\nolimits}
\def\CH{\mathop{\mathrm{CH}}\nolimits}
\def\proetale{{pro\text{-}\acute{e}tale}}
\def\etale{{\acute{e}tale}}
\def\QCoh{\mathit{QCoh}}
\def\Ker{\mathop{\mathrm{Ker}}}
\def\Im{\mathop{\mathrm{Im}}}
\def\Coker{\mathop{\mathrm{Coker}}}
\def\Coim{\mathop{\mathrm{Coim}}}

% Boxtimes
%
\DeclareMathSymbol{\boxtimes}{\mathbin}{AMSa}{"02}

%
% Macros for moduli stacks/spaces
%
\def\QCohstack{\mathcal{QC}\!\mathit{oh}}
\def\Cohstack{\mathcal{C}\!\mathit{oh}}
\def\Spacesstack{\mathcal{S}\!\mathit{paces}}
\def\Quotfunctor{\mathrm{Quot}}
\def\Hilbfunctor{\mathrm{Hilb}}
\def\Curvesstack{\mathcal{C}\!\mathit{urves}}
\def\Polarizedstack{\mathcal{P}\!\mathit{olarized}}
\def\Complexesstack{\mathcal{C}\!\mathit{omplexes}}
% \Pic is the operator that assigns to X its picard group, usage \Pic(X)
% \Picardstack_{X/B} denotes the Picard stack of X over B
% \Picardfunctor_{X/B} denotes the Picard functor of X over B
\def\Pic{\mathop{\mathrm{Pic}}\nolimits}
\def\Picardstack{\mathcal{P}\!\mathit{ic}}
\def\Picardfunctor{\mathrm{Pic}}
\def\Deformationcategory{\mathcal{D}\!\mathit{ef}}
\newenvironment{reference}{}{}
\newenvironment{slogan}{}{}
\newenvironment{history}{}{}
\usepackage{cleveref}
\crefname{originaltheorem}{original theorem}{original theorems}
\Crefname{originaltheorem}{Original theorem}{Original theorems}
\makeatletter
\newlabel{flat-theorem-flatten-module}{{0815}{source}{}{sourcepointer-1}{}}
\newlabel{flat-theorem-flatten-module@cref}{{[originaltheorem][0][]0815}{[1][0][]source}{}{}{}}
\makeatother
\crefname{originaltheorem}{original theorem}{original theorems}
\Crefname{originaltheorem}{Original theorem}{Original theorems}
\makeatletter
\newlabel{lemma-push-ideal}{{087B}{source}{}{sourcepointer-2}{}}
\newlabel{lemma-push-ideal@cref}{{[originaltheorem][0][]087B}{[1][0][]source}{}{}{}}
\makeatother
\crefname{originaltheorem}{original theorem}{original theorems}
\Crefname{originaltheorem}{Original theorem}{Original theorems}
\makeatletter
\newlabel{spaces-divisors-lemma-blowing-up-two-ideals}{{085Y}{source}{}{sourcepointer-3}{}}
\newlabel{spaces-divisors-lemma-blowing-up-two-ideals@cref}{{[originaltheorem][0][]085Y}{[1][0][]source}{}{}{}}
\makeatother
\crefname{originaltheorem}{original theorem}{original theorems}
\Crefname{originaltheorem}{Original theorem}{Original theorems}
\makeatletter
\newlabel{spaces-divisors-lemma-flat-base-change-blowing-up}{{085S}{source}{}{sourcepointer-4}{}}
\newlabel{spaces-divisors-lemma-flat-base-change-blowing-up@cref}{{[originaltheorem][0][]085S}{[1][0][]source}{}{}{}}
\makeatother
\crefname{originaltheorem}{original theorem}{original theorems}
\Crefname{originaltheorem}{Original theorem}{Original theorems}
\makeatletter
\newlabel{spaces-divisors-lemma-strict-transform-composition-blowups}{{0868}{source}{}{sourcepointer-5}{}}
\newlabel{spaces-divisors-lemma-strict-transform-composition-blowups@cref}{{[originaltheorem][0][]0868}{[1][0][]source}{}{}{}}
\makeatother
\crefname{originaltheorem}{original theorem}{original theorems}
\Crefname{originaltheorem}{Original theorem}{Original theorems}
\makeatletter
\newlabel{spaces-divisors-lemma-strict-transform-flat}{{0865}{source}{}{sourcepointer-6}{}}
\newlabel{spaces-divisors-lemma-strict-transform-flat@cref}{{[originaltheorem][0][]0865}{[1][0][]source}{}{}{}}
\makeatother
\crefname{originaltheorem}{original theorem}{original theorems}
\Crefname{originaltheorem}{Original theorem}{Original theorems}
\makeatletter
\newlabel{spaces-divisors-lemma-strict-transform-local}{{0863}{source}{}{sourcepointer-7}{}}
\newlabel{spaces-divisors-lemma-strict-transform-local@cref}{{[originaltheorem][0][]0863}{[1][0][]source}{}{}{}}
\makeatother
\title{Raynaud-Gruson flattening of a finite type module on a space by a U-admissible blowup}
\author{The Stacks Project Authors}
\date{stacks-project:spaces-more-morphisms.tex:local-snapshot-20261011; GNU FDL 1.2 or later}
\begin{document}
\maketitle
\noindent\textit{Editorial note.} Statement span is lines 8777-8794 of sources/stacks-spaces-more-morphisms/pinned-original-source.tex; proof span is lines 8796-8846, the complete author proof (49 body lines).\par\medskip
\noindent\textit{Editorial note.} Navigation stubs: 7 cited-but-not-included labels are supplied as non-load-bearing pointers with their pinned tags where the snapshot tag index has them.\par\medskip
\noindent\textit{Editorial note.} extra\_packages is empty on purpose: the retained preamble already loads xy with options and a second declaration clashes in the assembler wrapper.\par\medskip
\noindent\textit{Editorial note.} Text is retained verbatim, including the author's own typos and placeholders.\par\medskip
\noindent\textit{Editorial note.} The author's first paragraph of the schemes-side argument (More on Flatness) is quoted inside the chapter proof as the driver; no bibliography reference is invoked.\par\medskip
\noindent\textit{Allowed standard prerequisites (editorial).}\begin{itemize}
\item Explicit prerequisite boundary. Etale descent of strict transforms and of flatness/finite presentation (Divisors on Spaces, strict-transform-local, strict-transform-flat; Properties of Spaces descent), the schemes Raynaud-Gruson flattening theorem (More on Flatness, theorem-flatten-module), and the ideal-pushdown lemma (lemma-push-ideal, same chapter).
\end{itemize}\medskip
\ifdefined\hypertarget\hypertarget{sourcepointer-1}{}\fi
\noindent\textit{Source reference.} Original theorem 0815: sources/stacks-spaces-more-morphisms/pinned-original-source.tex, source lines 8846--8846. Cited by the retained statement or proof but not part of the collected spans; supplied only as a navigation pointer. The argument does not depend on its content.\par\medskip
\ifdefined\hypertarget\hypertarget{sourcepointer-2}{}\fi
\noindent\textit{Source reference.} Original theorem 087B: sources/stacks-spaces-more-morphisms/pinned-original-source.tex, source lines 8698--8712. Cited by the retained statement or proof but not part of the collected spans; supplied only as a navigation pointer. The argument does not depend on its content.\par\medskip
\ifdefined\hypertarget\hypertarget{sourcepointer-3}{}\fi
\noindent\textit{Source reference.} Original theorem 085Y: sources/stacks-spaces-more-morphisms/pinned-original-source.tex, source lines 8846--8846. Cited by the retained statement or proof but not part of the collected spans; supplied only as a navigation pointer. The argument does not depend on its content.\par\medskip
\ifdefined\hypertarget\hypertarget{sourcepointer-4}{}\fi
\noindent\textit{Source reference.} Original theorem 085S: sources/stacks-spaces-more-morphisms/pinned-original-source.tex, source lines 8846--8846. Cited by the retained statement or proof but not part of the collected spans; supplied only as a navigation pointer. The argument does not depend on its content.\par\medskip
\ifdefined\hypertarget\hypertarget{sourcepointer-5}{}\fi
\noindent\textit{Source reference.} Original theorem 0868: sources/stacks-spaces-more-morphisms/pinned-original-source.tex, source lines 8846--8846. Cited by the retained statement or proof but not part of the collected spans; supplied only as a navigation pointer. The argument does not depend on its content.\par\medskip
\ifdefined\hypertarget\hypertarget{sourcepointer-6}{}\fi
\noindent\textit{Source reference.} Original theorem 0865: sources/stacks-spaces-more-morphisms/pinned-original-source.tex, source lines 8846--8846. Cited by the retained statement or proof but not part of the collected spans; supplied only as a navigation pointer. The argument does not depend on its content.\par\medskip
\ifdefined\hypertarget\hypertarget{sourcepointer-7}{}\fi
\noindent\textit{Source reference.} Original theorem 0863: sources/stacks-spaces-more-morphisms/pinned-original-source.tex, source lines 8846--8846. Cited by the retained statement or proof but not part of the collected spans; supplied only as a navigation pointer. The argument does not depend on its content.\par\medskip
\begin{theorem}
\label{theorem-flatten-module}
Let $S$ be a scheme. Let $B$ be a quasi-compact and quasi-separated
algebraic space over $S$. Let $X$ be an algebraic space over $B$.
Let $\mathcal{F}$ be a quasi-coherent module on $X$.
Let $U \subset B$ be a quasi-compact open subspace. Assume
\begin{enumerate}
\item $X$ is quasi-compact,
\item $X$ is locally of finite presentation over $B$,
\item $\mathcal{F}$ is a module of finite type,
\item $\mathcal{F}_U$ is of finite presentation, and
\item $\mathcal{F}_U$ is flat over $U$.
\end{enumerate}
Then there exists a $U$-admissible blowup $B' \to B$ such that the
strict transform $\mathcal{F}'$ of $\mathcal{F}$ is an
$\mathcal{O}_{X \times_B B'}$-module of finite presentation and
flat over $B'$.
\end{theorem}

\begin{proof}
Choose an affine scheme $V$ and a surjective \'etale morphism $V \to X$.
Because strict transform commutes with \'etale localization
(Divisors on Spaces, Lemma \ref{spaces-divisors-lemma-strict-transform-local})
it suffices to prove the result with $X$ replaced by $V$. Hence we
may assume that $X \to B$ is representable (in addition to the
hypotheses of the lemma).

\medskip\noindent
Assume that $X \to B$ is representable. Choose an affine scheme $W$ and a
surjective \'etale morphism $\varphi : W \to B$. Note that $X \times_B W$
is a scheme. By the case of schemes
(More on Flatness, Theorem \ref{flat-theorem-flatten-module})
we can find a finite type quasi-coherent sheaf of ideals
$\mathcal{I} \subset \mathcal{O}_W$ such that
(a) $|V(\mathcal{I})| \cap |\varphi^{-1}(U)| = \emptyset$ and (b)
the strict transform of $\mathcal{F}|_{X \times_B W}$ with respect
to the blowing up $W' \to W$ in $\mathcal{I}$ becomes flat over $W'$
and is a module of finite presentation. Choose a finite type sheaf of ideals
$\mathcal{J} \subset \mathcal{O}_B$ as in Lemma \ref{lemma-push-ideal}.
Let $B' \to B$ be the blowing up of $\mathcal{J}$. We claim that this
blowup works. Namely, it is clear that $B' \to B$ is $U$-admissible
by our choice of ideal $\mathcal{J}$. Moreover, the base change
$B' \times_B W \to W$ is the blowup of $W$ in
$\varphi^{-1}\mathcal{J} = \mathcal{I}\mathcal{I}'$
(compatibility of blowup with flat base change, see
Divisors on Spaces, Lemma
\ref{spaces-divisors-lemma-flat-base-change-blowing-up}).
Hence there is a factorization
$$
W \times_B B' \to W' \to W
$$
where the first morphism is a blowup as well, see
Divisors on Spaces, Lemma \ref{spaces-divisors-lemma-blowing-up-two-ideals}).
The restriction of $\mathcal{F}'$ (which lives on $B' \times_B X$)
to $W \times_B B' \times_B X$ is the strict transform of
$\mathcal{F}|_{X \times_B W}$
(Divisors on Spaces, Lemma \ref{spaces-divisors-lemma-strict-transform-local})
and hence is the twice repeated strict transform of
$\mathcal{F}|_{X \times_B W}$ by the two blowups displayed above
(Divisors on Spaces, Lemma
\ref{spaces-divisors-lemma-strict-transform-composition-blowups}).
After the first blowup our sheaf is already flat over
the base and of finite presentation (by construction). Whence this holds 
after the second strict transform as well (since this is a
pullback by Divisors on Spaces, Lemma
\ref{spaces-divisors-lemma-strict-transform-flat}).
Thus we see that the restriction of $\mathcal{F}'$
to an \'etale cover of $B' \times_B X$ has the desired properties
and the theorem is proved.
\end{proof}

\end{document}
